\uuid{GrBU}
\exo7id{482}
\titre{exo7 482}
\auteur{cousquer}
\organisation{exo7}
\datecreate{2003-10-01}
\isIndication{false}
\isCorrection{false}
\chapitre{Propriétés de R}
\sousChapitre{Maximum, minimum, borne supérieure}
\cpge{mpsi-mp2i.nombres-reels-suites.b, pcsi.nombres-reels-suites.a, ptsi.nombres-reels-suites.a, mpsi-mp2i.fonctions-valeurs-reelles-complexes.a, pcsi.techniques-calcul-fonctions.a, ptsi.techniques-calcul-fonctions.a, bcpst1.nombres-reels, bcpst1.fonctions-reelles-usuelles.a}
\module{Analyse}
\niveau{L1}
\difficulte{}

\contenu{
\question{
Soit $A= \mathbb{Q}\cap\mathopen]0,1 \mathclose[$ et $a,b\in\mathbb{R}_+$.
On considère les applications suivantes de $A$ dans $\mathbb{R}_+$ :
$$ f \colon\frac{p}{q} \mapsto\frac{q-p}{q+p}\;; \qquad
g \colon\frac{p}{q} \mapsto\frac{aq+bp}{p+q}$$
Déterminer la borne supérieure et la borne inférieure de $f(A)$
et de $g(A)$.
}
}
